\chapter{Modül Detayları}

% ==========================================
% BÖLÜM 4.1: THISWORKBOOK (ÇALIŞMA KİTABI ETKİLEŞİM KATMANI)
% ==========================================

\section{\texttt{ThisWorkbook}}

\texttt{ThisWorkbook}, Excel VBA mimarisinde çalışma kitabının yaşam döngüsünü (\textit{Lifecycle}) ve global olayları (\textit{Workbook-Level Events}) dinleyen yerleşik sınıf modülüdür. Projede uygulamanın otomatik ilklendirilmesini (\textit{Initialization}) ve bellek değişkenlerinin hazırlanmasını sağlar.

Excel açıldığında grafik motoru, OLE bileşenleri ve Active Document yapısı henüz yüklenmemiş olabilir. Bu süreçte doğrudan makro çalıştırmak veya arayüze OLE nesneleri çizmek; grafik kaymalarına ve Korumalı Görünüm (\textit{Protected View}) modunda uygulamanın aniden çökmesine neden olur.

\texttt{ThisWorkbook} modülü, söz konusu kararsızlığı önlemek adına **Asenkron Gecikmeli Başlatma (Asynchronous Delayed Execution)** mimari desenini uygular:

\begin{itemize}\setlength{\itemsep}{0pt}\setlength{\parskip}{1pt}
    \item \textbf{Yerleşik Olay Yakalama (\texttt{Workbook\_Open}):} Otomatik tetiklenen giriş noktasıdır (\textit{Entry Point}).
    \item \textbf{İşlem Oluğu Sıfırlama (\texttt{DoEvents}):} Bekleyen tüm grafik yenileme ve ekran çizim olaylarını zorunlu kılar.
    \item \textbf{Zamanlanmış Tetikleme (\texttt{Application.OnTime}):} Arayüz kurma fonksiyonu olan \texttt{InitializeAppUI} prosedürünü 1 saniyelik mikro gecikmeyle Excel yürütme kuyruğuna kaydeder.
\end{itemize}

\vspace{0.15cm}

\begin{infobox}{\faClock\ Asenkron Başlatmanın Sistem Performansına Etkisi}
\texttt{Application.OnTime Now + TimeValue("00:00:01"), "InitializeAppUI"} kullanımı, dosya açılışında Excel UI motorunun tamamen stabil hale gelmesini bekler. Bu sayede servisler ve dinamik butonlar çökme riski olmadan sıfır hatayla ilklendirilir.
\end{infobox}

\vspace{0.15cm}

Aşağıdaki kod bloğu, projedeki \texttt{ThisWorkbook} sınıf modülünün tam yapısını içermektedir:

\begin{lstlisting}[caption={ThisWorkbook Sınıf Modülü Kod Yapısı}]
'*****************************************************************************
' Class: ThisWorkbook
' Purpose: Standard class module that handles workbook-level events.
'*****************************************************************************

Private Sub Workbook_Open()
    DoEvents
    Application.OnTime Now + TimeValue("00:00:01"), "InitializeAppUI"
End Sub
\end{lstlisting}

\newpage

% ==========================================
% BÖLÜM 4.2: CLSDYNAMICBUTTON (DİNAMİK NESNE TABANLI OLAY KONTROLÜ)
% ==========================================

\section{\texttt{clsDynamicButton}}

\texttt{clsDynamicButton}, çalışma zamanında (\textit{Runtime}) üretilen butonların tıklanma olaylarını (\texttt{Click Event}) Nesne Tabanlı Programlama (OOP) prensipleri doğrultusunda dinamik olarak yakalayan özel bir Sınıf Modülüdür (\textit{Class Module}).

UserForm kullanılmayan mimarilerde arayüz elemanları doğrudan \texttt{Worksheet} üzerine çizilir. Ancak dinamik eklenen bir \texttt{MSForms.CommandButton} nesnesinin statik bir tıklama prosedürü bulunmaz. Bu kısıtlamayı aşmak adına projedeki her bir dinamik buton, \texttt{clsDynamicButton} sınıfından türetilen bir nesneye bağlanır ve bellek ömrünü uzatmak için global bir koleksiyonda saklanır.

Sınıf içerisindeki \texttt{WithEvents} bildirimi, buton etkileşimlerini anlık dinler. Tıklama anında izlenen işlem adımları:

\begin{itemize}\setlength{\itemsep}{0pt}\setlength{\parskip}{1pt}
    \item \textbf{Referans ve Kapsayıcı Tespiti:} Tıklanan buton örneği yakalanır ve üst paneli (\texttt{MSForms.Frame}) doğrulanır.
    \item \textbf{Olay Delegasyonu:} Panelin adına göre (\texttt{pan\_Login}, \texttt{pan\_OSM} vb.) tetiklenecek yönlendirme \texttt{mod\_UI\_Events} modülündeki handler fonksiyonlarına aktarılır.
\end{itemize}

\vspace{0.15cm}

\begin{warningbox}{\faExclamationTriangle\ Bellek Yönetimi ve Nesne Yaşam Ömrü (Garbage Collection)}
VBA bellek yöneticisi, referansı tutulmayan sınıf örneklerini anında siler. Dinamik butonların tıklama olaylarının aktif kalması için \texttt{clsDynamicButton} örneklerinin \texttt{ButtonCollection} koleksiyonuna \texttt{Add} metoduyla eklenmesi kritik önem taşır.
\end{warningbox}

\vspace{0.15cm}

Aşağıdaki kod bloğu, dinamik buton olaylarını dinleyen ve ilgili UI panellerine yönlendiren sınıf yapısını özetlemektedir:

\begin{lstlisting}[caption={clsDynamicButton Sınıf Modülü Kod Yapısı}]
'*****************************************************************************
' Class: clsDynamicButton
' Purpose: This class is used to handle events for dynamically created buttons.
'*****************************************************************************

Public WithEvents DynamicButton As MSForms.CommandButton

Private Sub DynamicButton_Click()
    Dim parentFrame As MSForms.Frame
    
    On Error Resume Next
    Set parentFrame = DynamicButton.Parent
    On Error GoTo 0

    If parentFrame Is Nothing Then Exit Sub

    Select Case parentFrame.Name
        Case "pan_Login":    mod_UI_Events.HandleLoginPanelEvents DynamicButton.Name, parentFrame
        Case "pan_OSM":      mod_UI_Events.HandleOSMPanelEvents DynamicButton.Name, parentFrame
        Case "pan_Open":     mod_UI_Events.HandleOpenPanelEvents DynamicButton.Name, parentFrame
        ' . . . (pan_Closed, pan_PreviousMonth, pan_SelectPreviousMonth, ....)
    End Select
End Sub
\end{lstlisting}